% SPDX-License-Identifier: Apache-2.0
% covers: HostBridge
\datasheet{HostBridge}{The host side of an exit: AXI from a host in, packets into the network out.}

\dsfacts{\code{lib/parts/src/bus/noc/bridge.rs}}{\code{txhdl\_parts::bus::noc::bridge::HostBridge}}%
{\code{//docs:noc}}

\subsection*{Function}

\code{HostBridge} stands where an AXI host's five channels would go to
a peripheral, and sends them across the network instead. It reads the
host tracker's \code{aw}, \code{ar} and \code{w}, and drives its
\code{b} and \code{r}. On the network side it sends request packets on
\code{req} into its node's exit, and takes response packets from
\code{rsp}. It addresses each request by a map of three ranges.

\subsection*{Parameters}

\begin{center}\footnotesize
\begin{tabular}{@{}lp{0.7\textwidth}@{}}
\toprule
Parameter & Meaning \\
\midrule
\code{X}, \code{Y} & The column and row of the bridge's node. The
bridge stamps them on every request as its source. \\
\code{XB}, \code{YB} & The widths of a column and a row coordinate, in
bits. \\
\code{A} & The address width of the AXI link. \\
\code{D} & The data width of the AXI link. \\
\code{S} & The strobe width, which is \code{D / 8}. \\
\code{I} & The identifier width of the AXI link. \\
\code{B0}, \code{M0}, \code{X0}, \code{Y0} & The first range. An
address whose bits under the mask \code{M0} equal \code{B0} goes to
the node at \code{X0}, \code{Y0}. \\
\code{B1}, \code{M1}, \code{X1}, \code{Y1} & The second range, in the
same way. The first range wins where both match. \\
\code{B2}, \code{M2}, \code{X2}, \code{Y2} & The third range, the
default route. The doc comment says to give it a mask of zero. The
bridge's \code{run} does not read \code{B2} or \code{M2}: it sends
every address the first two ranges miss to \code{X2}, \code{Y2}. \\
\bottomrule
\end{tabular}
\end{center}

\subsection*{Address map}

The tables below are for \code{soc::Bridge<0, 0>}, the bridge of the
core at 0,0 in \code{soc/src/lib.rs}. That type is
\code{HostBridge<0, 0, 2, 2, 32, 32, 4, 2, 0x3000, 0xffff\_f000, 1, 1,
0x1000, 0xffff\_f000, 0, 1, 0, 0, 0, 1>}. Both hosts of \code{//soc}
use this map.

\begin{center}\footnotesize
\begin{tabular}{@{}llllp{0.4\textwidth}@{}}
\toprule
Range & Base & Mask & Node & What is there \\
\midrule
0 & \code{0x3000} & \code{0xffff\_f000} & 1,1 & The serial port. \\
1 & \code{0x1000} & \code{0xffff\_f000} & 0,1 & The program's data, in
the memory. \\
2 & \code{0} & \code{0} & 0,1 & Everything else, the display list and
the framebuffer included, in the memory. \\
\bottomrule
\end{tabular}
\end{center}

The masks reach the whole address. \code{//docs:noc} states why: with
a mask of \code{0xf000}, the framebuffer's rows at \code{0x13000} and
above went to the serial port.

\dstables{HostBridge}

\subsection*{Behaviour}

\begin{itemize}
\item The bridge holds nothing between cycles. A write is one packet,
so it has no burst to keep track of.
\item A write goes when \code{aw} offers an address phase, \code{w}
offers a beat, and \code{req} has room. The bridge takes both and sends
one packet with \code{chan} \code{W}. The packet holds the address
phase from \code{aw}, the data and strobe from \code{w}, and
\code{last} high.
\item A read goes when no write goes in that cycle, \code{ar} offers an
address phase, and \code{req} has room. The bridge takes it and sends
one packet with \code{chan} \code{Ar}.
\item The bridge sends at most one request packet per cycle, and a
write before a read.
\item Every request packet has the destination the map gives for its
address, the source \code{X}, \code{Y}, and the host's own identifier.
\item A response packet with \code{chan} \code{B} goes to \code{b} as
its identifier and response, when \code{b} is ready. Any other response
packet goes to \code{r} as its identifier, data, response and
\code{last}, when \code{r} is ready.
\end{itemize}

\subsection*{Verification}

\begin{itemize}
\item Three tests in \code{lib/parts/src/bus/noc.rs}, in
\code{//lib/parts:parts\_test}, put host bridges on a two by two
lattice with 16-bit addresses. Their map sends \code{0x1000} and
\code{0x2000}, each under the mask \code{0xf000}, and the default, all
to the node at 1,1.
\item \code{a\_burst\_crosses\_the\_lattice\_and\_its\_answer\_comes\_back}
has one host, at 0,0.
\item \code{two\_hosts\_share\_a\_memory\_and\_each\_answer\_goes\_home}
has hosts at 0,0 and 1,0 on one memory.
\item \code{four\_peripherals\_share\_one\_node} has hosts at 0,0 and
1,0, and four memories behind a router at 1,1.
\item \code{a\_core\_draws\_a\_solid\_and\_a\_gpu\_fills\_it}, in
\code{//soc:demo\_test}, runs \code{soc::Bridge} at 0,0 for the core
and at 1,0 for the rasteriser.
\item No waveform in \code{docs/BUILD.bazel} lowers the bridge, so its
netlist has no co-simulation under nvc or Verilator.
\code{//docs:showcase} lists the network's netlists as not yet
simulated.
\end{itemize}

\subsection*{Used by}

\code{soc::Bridge} in \code{soc/src/lib.rs}, for the core at 0,0 and
the rasteriser at 1,0. The tests of the network in
\code{lib/parts/src/bus/noc.rs}.

\subsection*{Limits}

A write of one beat is one packet. A burst of more than one beat does
not cross the network: it would need a virtual channel per source, or
reassembly at the peripheral, and the bridge has neither. The map has
three ranges. The last one takes every address the others miss. So no
address is unknown to the network, and neither bridge has to answer for
one.
